%----------------------------------------------------------------------
% 參數設定
% 這邊就是一些等一下模板在生出像是封面這些東西的時候，會需要用到的參數
%----------------------------------------------------------------------

% 中英文論文題目
\zhTitle{中文論文題目}
\enTitle{English Thesis Title}

% 中英文關鍵字
%       - 依據學校規定
%           關鍵詞 5-7 個，應附於摘要內
\zhKeywords{\textnormal{邊緣電漿子、過渡金屬二硫族化合物（TMDs）、時域有限差分法（FDTD）、色散關係、動量匹配}}
\enKeywords{Keyword1, Keyword2, Keyword3, Keyword4, Keyword5}

% 研究生中英文姓名
%       - 依據學校提供的範例
%           英文姓名應寫「姓, 名」，例：Wang, Jing
\zhStudentName{黃譯民}
\enStudentFirstName{Yi-Min}
\enStudentLastName{Huang}

% 指導教授中英文姓名
%       - 依據學校提供的範例
%           英文姓名應寫「姓, 名」，例：Wang, Jing
% 指導教授中英文姓名
\zhAdvisorName{吳致盛\ 博士、鄭舜仁\ 博士}
\enAdvisorFirstName{Chih-Sheng , Wu and Shun-Jen ,  Cheng}
\enAdvisorLastName{\newcommand{\gobble}[1]{}\gobble}
% 中英文學校名稱
\zhUnivName{國立陽明交通大學}
\enUnivName{National Yang Ming Chiao Tung University}

% 中英文學院名稱
\zhCollegeName{理學院}
\enCollegeName{}

% 中英文研究所名稱
\zhInstName{物理所}
\enInstName{IOP}

% 英文學位名稱
%       - 書名頁要用的
\enDegree{Master of Science}

% 英文領域名稱
%       - 書名頁要用的
%           出現在最下面的那一坨英文的最後一段文字
%       - 學校公版和系上的範本都有提供這個欄位
%           但是資科工所上傳審查的時候又說不用
%           所以如果不需要的話，可以留空就不會印東西出來。
\enField{Computer Science}

% 英文地點名稱
%       - 書名頁要用的
%           依據學校提供的範例，Taiwan, Republic of China
\enLocation{Taiwan, Republic of China}

% 論文中英文日期
\zhDegreeYear{一一○}
\zhDegreeMonth{七}
\enDegreeYear{2026}
\enDegreeMonth{July}

% Watermark of thesis (論文浮水印)
\watermark{Figures/watermark.png}
% ====================================================
% 強行拔除全本論文方程式編號的「式」字
% ====================================================
\usepackage{etoolbox}
\AtBeginDocument{
	\makeatletter
	\def\tagform@#1{\maketag@@@{(#1)}}
	\makeatother
}
